%%%%%%%%%%%%%%%%%%%%%%%%%%%%%%%%%%%%%%%%%%%%%%%%%%%%%%%%%%%%%%%%%%%%%%%%%%%%%%%%
%   Copyright 2009, Oracle and/or its affiliates.
%   All rights reserved.
%
%
%   Use is subject to license terms.
%
%   This distribution may include materials developed by third parties.
%
%%%%%%%%%%%%%%%%%%%%%%%%%%%%%%%%%%%%%%%%%%%%%%%%%%%%%%%%%%%%%%%%%%%%%%%%%%%%%%%%

\section{Proof of Overloading Resolution for Functions}
\seclabel{proof-overloading-resolution}

This section proves that the overloading rules for
function declarations are sufficient to guarantee no undefined nor
ambiguous call at run time (the case for methods is analogous).

The proof covers declarations whose parameter types and return types have no
static parameters, and calls to which declarations are applicable without
coercion.
It assumes that the declarations visible at the call satisfy the overloading
rules, so that no two of them have the same parameter type; that each of them has
a body and that the bodies are well typed; that the subtype relation is sound,
so that the dynamic type of an argument is a subtype of its static type, and
transitive; and that every value of a trait with a \KWD{comprises} clause is a
value of one of its listed types (\secref{trait-decls}).
The compiled type checker does not yet enforce the last assumption across
components: a type declared in a later component below a trait with static
parameters that stands between a trait with a \KWD{comprises} clause and its
listed types escapes its check (row~487 of the revival's gap ledger), and where
the assumption fails the uniqueness proved below is not guaranteed.
The proof does not extend to declarations with static parameters, whose
instantiation keeps obligations of its own (\secref{resolving-overloading},
\chapref{type-inference}), nor to coercion (\secref{proof-coercion-resolution}),
and it is not a proof of the soundness of the whole type system.

Consider a static function call $\f(\As)$ at some program point $Z$
and its corresponding dynamic function call $\f(\Xs)$.  Let $\Delta$
be the set of parameter types of function declarations of $\f$ that
are visible at $Z$ and applicable to the dynamic call $\f(\Xs)$.  Let
$\Sigma$ be the set of parameter types of function declarations of
$\f$ that are visible at $Z$ and applicable to the static call
$\f(\As)$.  Moreover, let $\sigma$ be the subset of $\Sigma$ for which
no type in $\Sigma$ is more specific and let $\delta$ be the
subset of $\Delta$ for which no type in $\Delta$ is more
specific:
\[
\begin{array}{lclcll}
\sigma & = & \set{S \in \Sigma & | & \lnot \exists S'\in\Sigma: S' \strictsubtype S & }
\\
\delta & = & \set{D \in \Delta & | & \lnot \exists D'\in\Delta: D' \strictsubtype D & }.
\end{array}
\]
For a declaration whose parameter type is $\Ps$, let $R_{\Ps}$ be its return
type.

From the Meet Rule (\secref{more-specific-rule}) the proof uses the following.
For any two declarations $\f(\Ps)$ and $\f(\Qs)$ visible at $Z$ such that
neither $\Ps$ nor $\Qs$ is a subtype of the other and $\Ps$ does not exclude
$\Qs$, there are declarations $\f(\Cs_1), \ldots, \f(\Cs_n)$ visible at $Z$
such that $\Cs_i \Ovrsubtype \Ps$ and $\Cs_i \Ovrsubtype \Qs$ for each $i$, and
every value of both $\Ps$ and $\Qs$ is a value of some $\Cs_i$: every value that
a program can make, not only the values of a particular call.
A declaration $\f(\Ps \inter \Qs)$ is the case of one such declaration.
This holds of every such pair of the declarations, the pairs among the
declarations $\f(\Cs_i)$ included.

Below we prove:
\[
\begin{array}{l}
|\Sigma| \neq 0 \ \text{implies} \ |\delta| = 1,
\\
\delta = \set{D} \ \text{implies} \ D \Ovrsubtype S \ \text{for every} \ S \in \Sigma, \ \text{and}
\\
\delta = \set{D} \ \text{implies} \ R_D \Ovrsubtype R_S \ \text{for every} \ S \in \sigma.
\end{array}
\]
Informally, if any declaration is applicable to a static call then there
exists a single most specific declaration that is applicable to the
corresponding dynamic call; it is the same as or more specific than every
declaration that is applicable to the static call; and its return type is a
subtype of the return type of every declaration applicable to the static call
that no other such declaration is more specific than, so that the value the
call returns belongs to the intersection of those return types.
The static call may have more than one most specific declaration: covering
gives a declaration that holds the value of the argument, not one whose
parameter type is a supertype of the static type of the argument.
\revision{revival-proof}{The Working Draft of February 2011 proved that
$|\sigma| = 1$ and $|\delta| = 1$ whenever $|\Sigma| \neq 0$, and that
$D \Ovrsubtype S$ where $\sigma = \set{S}$ and $\delta = \set{D}$, from a Meet
Rule that asks for a declaration on the intersection itself.
With the case of the Meet Rule in which declarations together cover the
intersection, the proof is revised: the lemmas of existence are kept, the
uniqueness of the most specific declaration for the static call is given up,
the uniqueness for the dynamic call is proved by a covering declaration, and
two steps replace the last theorem, the refinement of every declaration
applicable to the static call and the result type of the call.
The scope and the assumptions stated above are new.}

\begin{lemma}
\label{lem:subset}
$\Sigma \subseteq \Delta$.
\end{lemma}

\begin{proof}
Notice that $\Xs \Ovrsubtype \As$ by type soundness.
If $\f(\Ps)$ is applicable to
the call $\f(\As)$ then $\As \Ovrsubtype \Ps$.  Notice that $\Xs
\Ovrsubtype \As$ implies $\Xs \Ovrsubtype \Ps$.  Therefore
$\f(\Ps)$ is applicable to the call $\f(\Xs)$.
\end{proof}

\begin{lemma}
\label{lem:dgesge-second}
If $|\Delta| \ge 1$ then $|\delta| \ge 1$.
Also, if $|\Sigma| \ge 1$ then $|\sigma| \ge 1$.
\end{lemma}

\begin{proof}
Follows from Lemma \ref{lem:nonempty} where $S$ is the set of all
types, $A$ is $\Delta$ and $\Sigma$ respectively,
and the relation $\strictsubtype$ is acyclic and irreflexive.
\end{proof}

\begin{lemma}
\label{lem:ge}
If $|\Sigma| \ge 1$ then $|\delta| \ge 1$.
\end{lemma}

\begin{proof}
Follows from Lemmas \ref{lem:subset} and \ref{lem:dgesge-second}.
\end{proof}

\begin{lemma}
\label{lem:le}
$|\delta| \le 1$.
\end{lemma}

\begin{proof}
For the purpose of contradiction suppose there are two declarations
$\f(\Ps)$ and $\f(\Qs)$ in $\delta$.  Since both are in $\delta$ and
$\Ps \neq \Qs$, neither $\Ps$ nor $\Qs$ is a subtype of the other.
Since both $\f(\Ps)$ and $\f(\Qs)$ are applicable without coercion to the call
$\f(\Xs)$, the value of the argument is a value of both $\Ps$ and $\Qs$, so
$\Ps$ does not exclude $\Qs$.
By the Meet Rule as the proof uses it, there is a declaration $\f(\Cs)$
visible at $Z$ with $\Cs \Ovrsubtype \Ps$ and $\Cs \Ovrsubtype \Qs$ of which the
value is a value, so that $\Xs \Ovrsubtype \Cs$ and $\Cs \in \Delta$.
Since $\Ps$ is not a subtype of $\Qs$, $\Cs \neq \Ps$, so
$\Cs \strictsubtype \Ps$, which contradicts $\Ps \in \delta$.
\end{proof}

The same argument does not bound $|\sigma|$: the declaration $\f(\Cs)$ holds the
value of the argument, but $\Cs$ need not be a supertype of $\As$.

\begin{theorem}
\label{thm:overloading-subtyping}
If $|\Sigma| \neq 0$ then $|\delta| = 1$.
\end{theorem}

\begin{proof}
Follows from Lemmas \ref{lem:ge} and \ref{lem:le}.
\end{proof}

\begin{theorem}
\label{thm:dynamic-subtype-static}
If $\delta = \set{D}$ then $D \Ovrsubtype S$ for every $S \in \Sigma$.
\end{theorem}

\begin{proof}
Let $S \in \Sigma$.  By Lemma \ref{lem:subset}, $S \in \Delta$.
If $S \strictsubtype D$ then $D \notin \delta$, a contradiction.
Suppose that neither $S$ nor $D$ is a subtype of the other.
The value of the argument is a value of both, so $S$ does not exclude $D$, and
by the Meet Rule as the proof uses it there is a declaration $\f(\Cs)$ visible
at $Z$ with $\Cs \Ovrsubtype S$, $\Cs \Ovrsubtype D$ and $\Xs \Ovrsubtype \Cs$.
Then $\Cs \in \Delta$, and $\Cs \neq D$ since $D$ is not a subtype of $S$, so
$\Cs \strictsubtype D$, a contradiction.
Therefore $D \Ovrsubtype S$.
\end{proof}

\begin{theorem}
\label{thm:result-type}
If $\delta = \set{D}$ then $R_D \Ovrsubtype R_S$ for every $S \in \sigma$, and a
value that the dynamic call returns belongs to the intersection of the types
$R_S$ for $S \in \sigma$.
\end{theorem}

\begin{proof}
Let $S \in \sigma$.  By Theorem \ref{thm:dynamic-subtype-static},
$D \Ovrsubtype S$.  If $D = S$ the claim is immediate.
Otherwise $D \strictsubtype S$, and the only overloading rule that accepts such
a pair is the Subtype Rule (\secref{subtype-rule}), which requires
$R_D \Ovrsubtype R_S$.
The declaration with parameter type $D$ is the one that runs, and its body is
well typed, so a value it returns belongs to $R_D$, and so to $R_S$ for every
$S \in \sigma$.
This proves neither that the call terminates nor that nothing else fails at
run time.
\end{proof}
